\begin{ouverture}[geometrie]{17}{Symétrie axiale}{Partie 3 · Espace et géométrie}
{Axe · Symétrique d’un point · Constructions · Propriétés}
\objectif{Connaître la définition du symétrique d’un point par rapport à une droite.}
\objectif{Construire des symétriques sur quadrillage, papier pointé et papier uni.}
\objectif{Utiliser les propriétés de la symétrie pour construire et justifier.}
\prerequis{Médiatrice ; perpendiculaires ; report de longueurs au compas.}
\end{ouverture}
\pagestyle{otMain}\essentiel\fontsize{11}{15}\selectfont
\phantomsection\label{cours-17-1}
\essSection{1}{Reconnaître un axe et des points associés}
\begin{definition}Le symétrique d’un point M par rapport à une droite (d) est le point M' tel que (d) soit la médiatrice du segment [MM']. Si M est sur (d), son symétrique est M lui-même.\end{definition}
\label{fig-cours-17-1-cours-sym-point}\figureGeo{cours-sym-point}

\begin{definition}Une figure a un axe de symétrie si, pliée le long de cet axe, ses deux parties se superposent exactement. Un rectangle a 2 axes, un carré 4, un triangle équilatéral 3 ; un cercle en a une infinité : toutes les droites qui passent par son centre.\end{definition}
\essSection{2}{Construire sur supports gradués}
\begin{savoirfaire}[Quadrillage et papier pointé]Pour un axe horizontal ou vertical, on compte les carreaux de M jusqu’à l’axe, perpendiculairement à l’axe, puis on reporte le même nombre de carreaux de l’autre côté. Pour un axe qui suit les diagonales du quadrillage, on avance en diagonale à travers les carreaux. Sur papier pointé, on compte les points de la même façon.\end{savoirfaire}
\clearpage
\phantomsection\label{cours-17-2}
\essSection{3}{Construire sur papier uni}
\begin{savoirfaire}[À l’équerre et au compas]On trace la perpendiculaire à (d) passant par M ; elle coupe (d) en H. Au compas, on reporte la longueur MH de l’autre côté de (d) : on obtient M'. Au compas seul : avec deux points A et B de (d), M' est le second point commun au cercle de centre A passant par M et au cercle de centre B passant par M.\end{savoirfaire}
\label{fig-cours-17-2-cours-sym-triangle}\figureGeo{cours-sym-triangle}

\essSection{4}{Utiliser les propriétés de la symétrie}
\begin{propriete}La symétrie axiale conserve les longueurs, les angles, l’alignement et les aires. L’image d’un segment est un segment de même longueur ; l’image d’une droite est une droite ; l’image d’un cercle est un cercle de même rayon.\end{propriete}
\begin{savoirfaire}[Construire l’image d’une figure]Pour construire l’image d’un polygone, on construit les images de ses sommets, puis on les relie dans le même ordre. L’image d’un cercle se construit avec l’image de son centre et le même rayon.\end{savoirfaire}

\clearpage\rubrique{Découverte}
\begin{blocLivre}{ex-17-1}\begin{exercice}[1]{Combien d’axes ?}
Donner le nombre d’axes de symétrie de chaque figure : rectangle, carré, triangle isocèle, losange, triangle quelconque.
\label{fig-ex-17-1}\figureGeo{sym-1}
\end{exercice}\end{blocLivre}
\begin{blocLivre}{ex-17-2}\begin{exercice}[1]{Axe vertical}
Reproduire la figure et le point A, puis construire leurs symétriques par rapport à l’axe vertical.
\label{fig-ex-17-2}\figureGeo{sym-2}
\end{exercice}\end{blocLivre}
\begin{blocLivre}{ex-17-3}\begin{exercice}[1]{Axe horizontal}
Construire le symétrique du drapeau et du point B par rapport à l’axe horizontal.
\label{fig-ex-17-3}\figureGeo{sym-3}
\end{exercice}\end{blocLivre}
\begin{blocLivre}{ex-17-4}\begin{exercice}[1]{Axe en diagonale}
L’axe suit la diagonale des carreaux. Construire le symétrique de la figure et du point C.
\label{fig-ex-17-4}\figureGeo{sym-4}
\end{exercice}\end{blocLivre}
\begin{blocLivre}{ex-17-5}\begin{exercice}[1]{Papier pointé}
Construire le symétrique de la figure par rapport à la droite (d), sur papier pointé.
\label{fig-ex-17-5}\figureGeo{sym-5}
\end{exercice}\end{blocLivre}

\rubrique{Entraînement}
\begin{blocLivre}{ex-17-6}\begin{exercice}[2]{Symétrique d’un point}
Reproduire la droite (d) et le point A sur papier uni, puis construire le symétrique $A'$ de A par rapport à (d), à l’équerre et au compas.
\label{fig-ex-17-6}\figureGeo{sym-6}
\end{exercice}\end{blocLivre}
\begin{blocLivre}{ex-17-7}\begin{exercice}[2]{Image d’un segment}
Le segment [AB] mesure 4 cm. Construire son symétrique $[A'B']$ par rapport à la droite (d). Quelle est la longueur $A'B'$ ?
\label{fig-ex-17-7}\figureGeo{sym-7}
\end{exercice}\end{blocLivre}
\begin{blocLivre}{ex-17-8}\begin{exercice}[2]{Image d’un cercle}
Construire le symétrique du cercle de centre O et de rayon $1{,}5$ cm par rapport à la droite (d).
\label{fig-ex-17-8}\figureGeo{sym-8}
\end{exercice}\end{blocLivre}
\begin{blocLivre}{ex-17-9}\begin{exercice}[2]{Compléter la figure}
La droite (d) est un axe de symétrie de la maison. Compléter le dessin.
\label{fig-ex-17-9}\figureGeo{sym-9}
\end{exercice}\end{blocLivre}
\begin{blocLivre}{ex-17-10}\begin{exercice}[2]{Retrouver l’axe}
Le triangle $A'B'C'$ est le symétrique du triangle ABC par rapport à une droite (d) effacée. Retrouver la droite (d).
\label{fig-ex-17-10}\figureGeo{sym-10}
\end{exercice}\end{blocLivre}

\rubrique{Problèmes}
\begin{blocLivre}{ex-17-11}\begin{exercice}[3]{Ce que la symétrie conserve}
Le triangle $A'B'C'$ est le symétrique de ABC. On sait que $AB=5$ cm, $\widehat{ABC}=40^\circ$ et que l’aire de ABC vaut 6 cm$^2$. Donner, sans mesurer, $A'B'$, la mesure de l’angle de sommet $B'$ et l’aire de $A'B'C'$.
\label{fig-ex-17-11}\figureGeo{sym-11}
\end{exercice}\end{blocLivre}
\begin{blocLivre}{ex-17-12}\begin{exercice}[3]{Les points de l’axe}
Le segment [MN] coupe la droite (d) en I. Construire son symétrique $[M'N']$. Par quel point passe-t-il forcément ? Pourquoi ?
\label{fig-ex-17-12}\figureGeo{sym-12}
\end{exercice}\end{blocLivre}
\begin{blocLivre}{ex-17-13}\begin{exercice}[3]{Le billard}
On veut envoyer la boule A sur la bande (d) pour qu’elle rebondisse vers la boule B. Construire le symétrique $A'$ de A par rapport à (d), puis le point I où la droite $(A'B)$ coupe (d). Tracer le trajet A-I-B.
\label{fig-ex-17-13}\figureGeo{sym-13}
\end{exercice}\end{blocLivre}
\begin{blocLivre}{ex-17-14}\begin{exercice}[3]{Lettres symétriques}
Pour chacune des lettres H, A, E, N et X, dessinées ici en bâtons, donner le nombre d’axes de symétrie.
\label{fig-ex-17-14}\figureGeo{sym-14}
\end{exercice}\end{blocLivre}
\begin{blocLivre}{ex-17-15}\begin{exercice}[3]{Un retournement}
Malo pense que la symétrie axiale fait seulement glisser une figure. Observer le drapeau et son image. A-t-il raison ?
\label{fig-ex-17-15}\figureGeo{sym-15}
\end{exercice}\end{blocLivre}

\clearpage
